\SourceLeaf{363}%
umbeschriebenen Kreises mit \(AB\), ferner eine Strecke \(=AC-AB\) mit \(CF\). Durch Schlüsse in der bei den alten Geometern üblichen Form leitet er die folgenden Proportionen ab:
\par
{\fontsize{10}{12}\selectfont
\[
\begin{aligned}
AC:AB&=5AB+2CF:3AB+CF=19AB+7CF:11AB+4CF\\
&=71AB+26CF:41AB+15CF=265AB+97CF:153AB+56CF\\
&=989AB+362CF:571AB+209CF
\end{aligned}
\]
}
Aus der vorletzten folgt dann leicht \(AC:AB>265:153\). Aus der letzten folgt bei Einführung einer Strecke \(BD=2AB-AC=AB-CF\)
\[
AC:AB=1351AB-362BD:780AB-209BD<1351:780
\]
Wir räumen gern ein, dass Herr M., der sich mit der bei den alten Geometern üblichen Form arithmetischer Schlüsse sehr vertraut gemacht hat, den Gedankengang des \textsc{Archimedes} tatsächlich erraten haben könnte. Umso weniger lässt sich darüber entscheiden, als solche Untersuchungen auf sehr vielfältige Weise angegangen werden können. Auch gibt es sonst Anzeichen dafür, dass der große Grieche manche keineswegs gewöhnlichen Wahrheiten und Kunstgriffe gekannt haben muss, selbst aus der höheren Arithmetik.

Eine Frage bleibt hier noch offen. Warum ist \textsc{Archimedes}, wenn er seine Näherungswerte methodisch gefunden hat. bei den größeren bis zum elften Glied gegangen, während er bei den kleineren nur bis zum achten ging? Man sollte meinen, dass er sich bei jenen mit dem neunten Glied \(\frac{362}{209}\) begnügt hätte, das auf jeden Fall ausgereicht hätte, um die untere Grenze \(3\frac{10}{71}\) zu bestimmen. Man könnte sich vielleicht zu dem Schluss verleiten lassen, dass \textsc{Archimedes} den Bruch \(\frac{1351}{780}\) doch durch eine Art glücklichen Zufall gefunden hat und ihm der einfachere \(\frac{362}{209}\) entgangen ist. Herr M. glaubt, \textsc{Archimedes} habe jenen Bruch gewählt, weil er der einfachste unter denen mit einem Zähler in der Größenordnung der Tausender sei, so wie er den Bruch \(\frac{265}{153}\) als den einfachsten aus der Größenordnung der Hunderter gewählt habe. Dieser Grund scheint uns jedoch nicht befriedigend. Wir halten es für wahrscheinlicher, dass er den Bruch \(\frac{1351}{780}\) vorzog, weil sich damit im weiteren Verlauf der Rechnung zufällig eine bequeme Vereinfachung ergab. So ließ sich beim 24-Eck eine äußerst nahe Grenze für das Verhältnis, das in unserer Ausdrucksweise \(1:\operatorname{cotang}7^0\,30'\) ist, sehr %
